% GENERATED FILE. Edit canonical lessons, then run npm run generate:books.
% canonical-source-hash: fnv1a64:a9334829396207b3
% canonical-lessons: MR-C395-upyukta, MR-C395-avghad, MR-C395-thet, MR-C395-apratyaksa, MR-C395-aspasta, MR-R395-review-389-395, MR-R395-review-final-395

\chapter{Useful or Not}
\label{ch:mr-useful-or-not}
\hlchaptermodality{395}

\begin{chapteropening}
\textbf{By the end of this chapter:} \emph{I can say whether something is useful, tricky, direct or unclear.}

Complete the last lesson of chapter 395: I can say whether something is useful, tricky, direct or unclear.
\end{chapteropening}

\section[\texorpdfstring{\mr{उपयुक्त} (upyukta)}{उपयुक्त (upyukta)}]{\mr{उपयुक्त} (upyukta) — useful}

\label{lesson:MR-C395-upyukta}



\begin{quote}
Before the new one: say the Marathi for good judgement, then the Marathi for to dawn on.
\end{quote}

\subsection*{You'll want to know: \mr{उपयुक्त}}
\textbf{\mr{उपयुक्त}} — \emph{upyukta} — \textquotedblleft{}useful\textquotedblright{}.

\begin{quote}
\mr{हे} \mr{पुस्तक} \mr{परीक्षेसाठी} \mr{खूप} \mr{उपयुक्त} \mr{आहे}.
\end{quote}

\begin{quote}
\emph{he pustak parīkṣesāṭhī khūp upyukta āhe.}
\end{quote}

\begin{quote}
This book is very useful for the exam.
\end{quote}

Invariable; X \mr{साठी} \mr{उपयुक्त}, useful for X.

\begin{cousinweb}
\textbf{\mr{उपयुक्त}} (\emph{upyukta}) is a tatsama, taken from Sanskrit \emph{upayukta}, \textquotedblleft{}fit, suitable, useful\textquotedblright{} (Monier-Williams).
\end{cousinweb}

\begin{grammarlens}[title={What you've built}]
One more word for saying what happens if.
\end{grammarlens}

\subsection*{Your turn}
\begin{itemize}
\raggedright
  \item \emph{Say it:} \emph{upyukta}
  \item \emph{Say it:} \emph{upyukta}, once more
  \item \emph{Say it:} \emph{he pustak parīkṣesāṭhī khūp upyukta āhe.} — This book is very useful for the exam.
  \item \emph{Recall:} say the Marathi for good judgement, then the Marathi for to dawn on, then say \emph{upyukta} again
\end{itemize}

\begin{tcolorbox}[breakable,colback=teal!4,colframe=teal!35!black,title={Before you move on}]
What does \mr{उपयुक्त} mean? (\textquotedblleft{}Useful\textquotedblright{}.) Say it once more.
\end{tcolorbox}
\section[\texorpdfstring{\mr{अवघड} (avghaḍ)}{अवघड (avghaḍ)}]{\mr{अवघड} (avghaḍ) — tricky; hard}

\label{lesson:MR-C395-avghad}



\begin{quote}
Before the new one: say the Marathi for to dawn on, then the Marathi for useful.
\end{quote}

\subsection*{You'll want to know: \mr{अवघड}}
\textbf{\mr{अवघड}} — \emph{avghaḍ} — \textquotedblleft{}tricky; hard\textquotedblright{}.

\begin{quote}
\mr{गणिताचा} \mr{पेपर} \mr{अवघड} \mr{होता}.
\end{quote}

\begin{quote}
\emph{gaṇitācā pepar avghaḍ hotā.}
\end{quote}

\begin{quote}
The maths paper was hard.
\end{quote}

Invariable; \mr{अवघड} \mr{जागी}, in an awkward spot.

\begin{cousinweb}
\textbf{\mr{अवघड}} (\emph{avghaḍ}) is a Marathi word whose origin is not settled.
\end{cousinweb}

\begin{grammarlens}[title={What you've built}]
One more word for saying what happens if.
\end{grammarlens}

\subsection*{Your turn}
\begin{itemize}
\raggedright
  \item \emph{Say it:} \emph{avghaḍ}
  \item \emph{Say it:} \emph{avghaḍ}, once more
  \item \emph{Say it:} \emph{gaṇitācā pepar avghaḍ hotā.} — The maths paper was hard.
  \item \emph{Recall:} say the Marathi for to dawn on, then the Marathi for useful, then say \emph{avghaḍ} again
\end{itemize}

\begin{tcolorbox}[breakable,colback=teal!4,colframe=teal!35!black,title={Before you move on}]
What does \mr{अवघड} mean? (\textquotedblleft{}Tricky; hard\textquotedblright{}.) Say it once more.
\end{tcolorbox}
\section[\texorpdfstring{\mr{थेट} (theṭ)}{थेट (theṭ)}]{\mr{थेट} (theṭ) — direct; straight}

\label{lesson:MR-C395-thet}



\begin{quote}
Before the new one: say the Marathi for useful, then the Marathi for tricky.
\end{quote}

\subsection*{You'll want to know: \mr{थेट}}
\textbf{\mr{थेट}} — \emph{theṭ} — \textquotedblleft{}direct; straight\textquotedblright{}.

\begin{quote}
\mr{ही} \mr{बस} \mr{थेट} \mr{पुण्याला} \mr{जाते}.
\end{quote}

\begin{quote}
\emph{hī bas theṭ puṇyālā jāte.}
\end{quote}

\begin{quote}
This bus goes straight to Pune.
\end{quote}

Also an adverb: \mr{थेट} \mr{सांगणे}, to say straight out.

\begin{cousinweb}
\textbf{\mr{थेट}} (\emph{theṭ}) is a Marathi word whose origin is not settled. CDIAL lists a Marathi \emph{theṭ}, \textquotedblleft{}genuine\textquotedblright{}, under Sanskrit \emph{sthēṣṭha}, \textquotedblleft{}most fixed, durable\textquotedblright{} (CDIAL 13778); whether the sense direct comes from it is not clear.
\end{cousinweb}

\begin{grammarlens}[title={What you've built}]
One more word for saying what happens if.
\end{grammarlens}

\subsection*{Your turn}
\begin{itemize}
\raggedright
  \item \emph{Say it:} \emph{theṭ}
  \item \emph{Say it:} \emph{theṭ}, once more
  \item \emph{Say it:} \emph{hī bas theṭ puṇyālā jāte.} — This bus goes straight to Pune.
  \item \emph{Recall:} say the Marathi for useful, then the Marathi for tricky, then say \emph{theṭ} again
\end{itemize}

\begin{tcolorbox}[breakable,colback=teal!4,colframe=teal!35!black,title={Before you move on}]
What does \mr{थेट} mean? (\textquotedblleft{}Direct; straight\textquotedblright{}.) Say it once more.
\end{tcolorbox}
\section[\texorpdfstring{\mr{अप्रत्यक्ष} (apratyakṣa)}{अप्रत्यक्ष (apratyakṣa)}]{\mr{अप्रत्यक्ष} (apratyakṣa) — indirect}

\label{lesson:MR-C395-apratyaksa}



\begin{quote}
Before the new one: say the Marathi for tricky, then the Marathi for direct.
\end{quote}

\subsection*{You'll want to know: \mr{अप्रत्यक्ष}}
\textbf{\mr{अप्रत्यक्ष}} — \emph{apratyakṣa} — \textquotedblleft{}indirect\textquotedblright{}.

\begin{quote}
\mr{त्याच्या} \mr{बोलण्यात} \mr{अप्रत्यक्ष} \mr{टीका} \mr{होती}.
\end{quote}

\begin{quote}
\emph{tyācyā bolaṇyāt apratyakṣa ṭīkā hotī.}
\end{quote}

\begin{quote}
There was indirect criticism in his words.
\end{quote}

\mr{अ} + \mr{प्रत्यक्ष}, not before the eyes.

\begin{cousinweb}
\textbf{\mr{अप्रत्यक्ष}} (\emph{apratyakṣa}) is a tatsama, taken from Sanskrit \emph{apratyakṣa}, \textquotedblleft{}not present to the eyes, invisible\textquotedblright{} (Monier-Williams).
\end{cousinweb}

\begin{grammarlens}[title={What you've built}]
One more word for saying what happens if.
\end{grammarlens}

\subsection*{Your turn}
\begin{itemize}
\raggedright
  \item \emph{Say it:} \emph{apratyakṣa}
  \item \emph{Say it:} \emph{apratyakṣa}, once more
  \item \emph{Say it:} \emph{tyācyā bolaṇyāt apratyakṣa ṭīkā hotī.} — There was indirect criticism in his words.
  \item \emph{Recall:} say the Marathi for tricky, then the Marathi for direct, then say \emph{apratyakṣa} again
\end{itemize}

\begin{tcolorbox}[breakable,colback=teal!4,colframe=teal!35!black,title={Before you move on}]
What does \mr{अप्रत्यक्ष} mean? (\textquotedblleft{}Indirect\textquotedblright{}.) Say it once more.
\end{tcolorbox}
\section[\texorpdfstring{\mr{अस्पष्ट} (aspaṣṭa)}{अस्पष्ट (aspaṣṭa)}]{\mr{अस्पष्ट} (aspaṣṭa) — unclear}

\label{lesson:MR-C395-aspasta}



\begin{quote}
Before the new one: say the Marathi for direct, then the Marathi for indirect.
\end{quote}

\subsection*{You'll want to know: \mr{अस्पष्ट}}
\textbf{\mr{अस्पष्ट}} — \emph{aspaṣṭa} — \textquotedblleft{}unclear\textquotedblright{}.

\begin{quote}
\mr{फोनवर} \mr{तुझा} \mr{आवाज} \mr{अस्पष्ट} \mr{येत} \mr{आहे}.
\end{quote}

\begin{quote}
\emph{phonvar tujhā āvāj aspaṣṭa yet āhe.}
\end{quote}

\begin{quote}
Your voice is coming through unclear on the phone.
\end{quote}

\mr{अ} + \mr{स्पष्ट}, not clear.

\begin{cousinweb}
\textbf{\mr{अस्पष्ट}} (\emph{aspaṣṭa}) is a tatsama, taken from Sanskrit \emph{aspaṣṭa}, \textquotedblleft{}not clear, indistinct\textquotedblright{} (Monier-Williams).
\end{cousinweb}

\begin{grammarlens}[title={What you've built}]
One more word for saying what happens if.
\end{grammarlens}

\subsection*{Your turn}
\begin{itemize}
\raggedright
  \item \emph{Say it:} \emph{aspaṣṭa}
  \item \emph{Say it:} \emph{aspaṣṭa}, once more
  \item \emph{Say it:} \emph{phonvar tujhā āvāj aspaṣṭa yet āhe.} — Your voice is coming through unclear on the phone.
  \item \emph{Recall:} say the Marathi for direct, then the Marathi for indirect, then say \emph{aspaṣṭa} again
\end{itemize}

\begin{tcolorbox}[breakable,colback=teal!4,colframe=teal!35!black,title={Before you move on}]
What does \mr{अस्पष्ट} mean? (\textquotedblleft{}Unclear\textquotedblright{}.) Say it once more.
\end{tcolorbox}
\section[(dialogue)]{Looking back on conditions, degrees and judgement}

\label{lesson:MR-R395-review-389-395}



\begin{quote}
Say the Marathi for partly and completely.
\end{quote}

\subsection*{Your turn}
Say these aloud:

\begin{itemize}
\raggedright
  \item partly, completely, merely, purely, usually
  \item in total, on the whole, a principle, particular, apart from
  \item to decline, to hold up, a percent, a proportion, a level
  \item a clause, to allocate, to pass over, to follow, to embrace
  \item an excuse, a pretext, to brush aside, to request, to follow up
  \item awareness, a realization, alertness, good judgement, to dawn on
  \item useful, tricky, direct, indirect, unclear
\end{itemize}

\begin{tcolorbox}[breakable,colback=teal!4,colframe=teal!35!black,title={Before you move on}]
Partly? (\textbf{\mr{अंशतः}}, \emph{aṁśataḥ}.) Completely? (\textbf{\mr{पूर्णपणे}}, \emph{pūrṇapaṇe}.) Merely? (\textbf{\mr{केवळ}}, \emph{kevaḷ}.) Purely? (\textbf{\mr{निव्वळ}}, \emph{nivvaḷ}.) Usually? (\textbf{\mr{सहसा}}, \emph{sahsā}.) In total? (\textbf{\mr{एकूण}}, \emph{ekūṇ}.) On the whole? (\textbf{\mr{एकंदरीत}}, \emph{ekandarīt}.) A principle? (\textbf{\mr{तत्त्व}}, \emph{tattva}.) Particular? (\textbf{\mr{विशेष}}, \emph{viśeṣ}.) Apart from? (\textbf{\mr{व्यतिरिक्त}}, \emph{vyatirikta}.) To decline? (\textbf{\mr{घटणे}}, \emph{ghaṭṇe}.) To hold up? (\textbf{\mr{टिकणे}}, \emph{ṭikṇe}.) A percent? (\textbf{\mr{टक्का}}, \emph{ṭakkā}.) A proportion? (\textbf{\mr{प्रमाण}}, \emph{pramāṇ}.) A level? (\textbf{\mr{स्तर}}, \emph{star}.) A clause? (\textbf{\mr{कलम}}, \emph{kalam}.) To allocate? (\textbf{\mr{वाटप} \mr{करणे}}, \emph{vāṭap karṇe}.) To pass over? (\textbf{\mr{डावलणे}}, \emph{ḍāvalṇe}.) To follow? (\textbf{\mr{अनुसरणे}}, \emph{anusarṇe}.) To embrace? (\textbf{\mr{अंगीकारणे}}, \emph{aṅgīkārṇe}.) An excuse? (\textbf{\mr{सबब}}, \emph{sabab}.) A pretext? (\textbf{\mr{बहाणा}}, \emph{bahāṇā}.) To brush aside? (\textbf{\mr{धुडकावणे}}, \emph{dhuḍkāvṇe}.) To request? (\textbf{\mr{विनंती} \mr{करणे}}, \emph{vinantī karṇe}.) To follow up? (\textbf{\mr{पाठपुरावा} \mr{करणे}}, \emph{pāṭhpurāvā karṇe}.) Awareness? (\textbf{\mr{भान}}, \emph{bhān}.) A realization? (\textbf{\mr{जाणीव}}, \emph{jāṇīv}.) Alertness? (\textbf{\mr{जागरूकता}}, \emph{jāgarūktā}.) Good judgement? (\textbf{\mr{विवेक}}, \emph{vivek}.) To dawn on? (\textbf{\mr{उमगणे}}, \emph{umagṇe}.) Useful? (\textbf{\mr{उपयुक्त}}, \emph{upyukta}.) Tricky? (\textbf{\mr{अवघड}}, \emph{avghaḍ}.) Direct? (\textbf{\mr{थेट}}, \emph{theṭ}.) Indirect? (\textbf{\mr{अप्रत्यक्ष}}, \emph{apratyakṣa}.) Unclear? (\textbf{\mr{अस्पष्ट}}, \emph{aspaṣṭa}.)
\end{tcolorbox}
\section[(dialogue)]{Useful or Not — once more}

\label{lesson:MR-R395-review-final-395}



\begin{quote}
Say the Marathi for useful and tricky.
\end{quote}

\subsection*{Your turn}
\begin{itemize}
\raggedright
  \item \emph{Say it:} useful, tricky, direct, indirect, unclear
\end{itemize}

\begin{tcolorbox}[breakable,colback=teal!4,colframe=teal!35!black,title={Before you move on}]
Useful? (\textbf{\mr{उपयुक्त}}, \emph{upyukta}.) Tricky? (\textbf{\mr{अवघड}}, \emph{avghaḍ}.) Direct? (\textbf{\mr{थेट}}, \emph{theṭ}.) Indirect? (\textbf{\mr{अप्रत्यक्ष}}, \emph{apratyakṣa}.) Unclear? (\textbf{\mr{अस्पष्ट}}, \emph{aspaṣṭa}.)
\end{tcolorbox}
